% Nouveau EGA V (cote 25): archive cote 25, site page 59 (PDF page 60).
% typescript. Draft compared: none.
% AI-transcribed by Claude Opus 5.5 (claude-opus-5-5, adaptive thinking, effort high), Message Batches API, batch ega25b (msgbatch_01A7n5YuXubUPNK1zsepi8RG), 2026-11-13; not yet checked by eye.
\scanpage{25/59}
\mspage{59}

Comme minimum dans ce sens, il faut  donner ici la reformulation de 9.5.
en termes de pinceaux linéaires  -  fait constamment utilisé dans les
applications géométriques.
\hand{1}0.2.  Par polarité, une droite $\underline{L}$  de $\check{P}$  correspond à une sous-variété
linéaire $L^{o}$ de codimension 2 dans P  (S quélconque).  Posons
\[ T = X\mathrm{x}_{P}L^{o} \]
Une autre façon  de décrire T est la suivante:  L  est défini par un
quotient loc libre de rang 2 de $\check{\underline{E}}$, ou ce qui revient au même, par un
sous-Module localement facteur direct  $\underline{F}$ de $\underline{E}$\xout{x}, partout de rang 2,
Considérons  l'homomorphisme composé
\[ \underline{F}_{X} \longrightarrow \underline{E}_{X} \longrightarrow \underline{O}_{X}(1) \quad , \]
alors  T n'est autre que  le schéma des zéros de cet homomorphisme compo-%
sé, où ce qui revient au même, est défini par l'Idéal \inserted[typed]{J} image de l'homomo%
phisme correspondant (obtenu en tordant par $\underline{O}_{X}(-1)$):
\[ \underline{F}_{X}(-1) \longrightarrow \underline{O}_{X} \quad . \]
\inserted[typed]{Supposons}
\xout{Supposonsxx Pour} que cet homomorphisme soit \uline{régulier}, ce qui signifie
que  si on écrit localement une base \inserted{tot. ordonnée} de $\underline{F}_{X}(-1)$, \deleted{leurs} \inserted{les} images dans $\underline{O}_{X}$\hand{\unsure{\^{}}}
forme\over{\unsure{nt}}{nt} une suite $\underline{O}_{X}$-régulière, \xout{x}condition qui ne dépend pas de la base
choisie et peut s'énoncer intrinsèquement aussi en disant que
$\underline{F}_{X}(-1)$\inserted{$\otimes\ \underline{O}_{X}/J$}$\to \underline{J}/\underline{J}^{2}$  est un isomorphisme, et  $V(\underline{J})=T \longrightarrow X$  est  une immersion
\margin{|} régulière  (NB il \xout{xxxxxx} faudrait quelque part dégager la situation
\margin{|} générale avec un hom  $\underline{G} \to \underline{O}_{X}$ , $\underline{G}$ localement libre sur préschéma X, par
\margin{|} exemple dans le par. des immersions régulières).  Alors on a la
\xout{Proposition} \uline{Théorème}  10.2.   Sous l'hypothèse précédente, le pinceau
linéaire $\underline{Y}_{L}$  muni de la projection canonique $\underline{Y}_{L} \longrightarrow X$  est  X-isomorphe
\inserted{de façon unique} au préschéma éclaté  de X au moyen de T.

\mspage{59}
